\scanpage{72}
2. Je suppose maintenant que $R$ est un esp. top \inserted{séparé} et je suppose

\textcircled{\unsure{7|6}} Si \inserted{$n\in\mathbb N$,} $X\in C_n$, $X$ est \uline{compact}.

Cela implique que si $X\in C_n$, $Y\in C_p$, alors tt $C_*$-morphisme de $X$ dans $Y$ est \uline{continu}, que les $C_*$-modèles sont de façon canonique munis d'une structure d'espace compact, \deleted{qui} \unsure{sorte qu'on trouve} un foncteur canonique
\[
(C_*\text{-modèles})\longrightarrow\text{espaces compacts}.
\]

Je suppose de plus

\textcircled{8} Si $n\in\mathbb N$, $X\in C_n$, et $x\in X$, alors \deleted{$\forall$ voisinage $V$ de \unsure{$x$}; les}
\deleted{$V\in C_n$ qui sont un voisinage de $x$ dans $X$, $\exists Y$}
les $Y\in C_n$ qui \inserted{$\subset X$} sont un voisinage de $x$ \inserted{dans $X$} forment un syst. fond. de voisinages de $x$ dans $X$.

[Il suffit que \deleted{\unsure{tous}} les intérieurs des $X\in C_n$ forment une base de la topologie de $R^n$]
\margin{C'est ce qu'on supposera}

Alors, les \inserted{$C_*$-}morphismes d'un $C_*$-modèle $X$ dans un autre $Y$ sont les sections d'un faisceau -- le faisceau des germes de $C_*$-morphismes de $X$ dans $Y$.
Idem pour les $C_*$-morphismes dans $R$ (ou dans un $R^I$).
La connaissance du faisceau des $C_*$-morphismes d'un modèle $X$ (donné comme espace cpt) dans $R$, \inserted{\unsure{connu pour ss-faisceau de celui des appl. cont. de $X$ dans $R$|contt pour ss-faisceau de celui des appl. cont. de $X$ dans $R$}}, redonne la \inserted{$C_*$-}structure du modèle.

$C_*$-espace : espace loc.$^{\mathrm t}$ compact, muni d'un sous-faisceau du faisceau des appl. continues à val. dans $R$, tel que tt pt ait un voisinage cpt qui, pour la structure induite, soit un $C_*$-modèle ---

\margin{73}
